
\section{Introduction}

Perception systems for open-vocabulary segmentation and OCR are still largely built around a modular encoder--decoder recipe: a vision backbone produces features, and a separate decoder or late-fusion module turns them into task outputs \citep{carion2020detr,chen2022pix2seq,kamath2021mdetr,kirillov2023segment,sam2}. This design has been effective, but it also encourages adding task-specific mechanisms (modality fusion, query matching, post-processing) that complicate scaling and limit effective feature fusion and learning between vision and text modality \citep{carion2025sam3,minderer2024scalingopenvocabularyobjectdetection,Liu2023GroundingDM,li2022grounded}.

This raises a fundamental question: \emph{do dense grounding systems actually need an encoder--decoder split to ``see'' and to ``predict''?} If not, what should replace the standard late-fusion interface so that language and visual features can interact from the first layer, especially for compositional prompts involving spatial constraints, and relations. What is the right output interface for dense perception, where the number of instances can range from zero to hundreds, without making decoding prohibitively expensive \citep{chen2022pix2seq,chen2022unified,kolesnikov2022uvim}? And finally, how should we benchmark progress once standard referring datasets saturate, in a way that isolates these compositional capabilities and stresses long-context crowded scenes?

In this report, we present Falcon Perception, a unified dense transformer model for dense vision-language perception. Given an image and a prompt, the model must decide \emph{whether the concept is present} and, if so, identify the referred instance(s) and predict pixel masks~\citep{carion2025sam3,kamath2021mdetr}. This setting covers open-vocabulary / promptable segmentation and language-grounded perception~\citep{kirillov2023segment,sam2,li2022grounded,Liu2023GroundingDM,minderer2024scalingopenvocabularyobjectdetection}. While standard approaches typically decouple visual feature extraction (via a vision encoder) from task-specific decoding (via a separate transformer decoder). Falcon Perception unifies these capabilities into a single dense stack. By processing image patches and text tokens in a shared parameter space from the first layer, we avoid late fusion bottlenecks and keep the system simpler for dense grounding.


\noindent Our architecture is built on the hypothesis that a single transformer, equipped with the right inductive biases, can simultaneously learn good visual representations and support autoregressive task generation. To this end, we introduce several key architectural and training innovations:

\begin{enumerate}
    \item \textbf{Unified Dense Transformer with Hybrid Attention Mask:} We replace the encoder--decoder split with a single Transformer using a hybrid attention mask: image tokens attend bidirectionally to build global context, while text and task tokens attend causally to the full visual prefix. This induces encoder-like behavior for vision and decoder-like behavior for language within the same weights.

    \item We introduce \textbf{PBench}, a benchmark for compositional prompts and dense long-context evaluation, enabling capability-level analysis beyond saturated referring benchmarks.


    \item \textbf{Chain-of-Perception:} To balance model expressivity with fast inference, we decompose each instance segmentation into a sequence: \texttt{<coord>} $\rightarrow$ \texttt{<size>} $\rightarrow$ \texttt{<seg>}. This ordering forces the model to resolve an instance's position and size as conditioning signal before producing its segmentation mask.
    
    \item \textbf{Specialized heads:} The backbone is shared, but decoding is not: each Chain-of-Perception token has its own lightweight head. This makes coordinate/size prediction structured (via Fourier features) and keeps segmentation fast by generating high-resolution masks in parallel (Section~\ref{sec:heads}).
    
\end{enumerate}

\noindent Finally, we evaluate Falcon Perception on both referring expression segmentation and open-vocabulary segmentation benchmarks, where it compares favorably to state-of-the-art systems. We also extend the same early-fusion recipe to text-heavy vision tasks with Falcon-OCR and show that a compact 300M model can reach strong OCR performance relative to much larger systems. Falcon Perception connects to these lines of work:

\noindent \textbf{Open-vocabulary segmentation and grounding:} The Segment Anything family introduced promptable segmentation with interactive refinement~\citep{kirillov2023segment,sam2}. SAM~3~\citep{carion2025sam3} formalizes \emph{promptable concept segmentation} and introduces SA-Co together with a calibration-aware evaluation protocol (pmF$_1$, IL\_MCC, and cgF$_1$). A key takeaway for open vocabulary is that recognition and localization should be disentangled: the model must learn to say ``absent'' reliably, not only draw masks when confident. In parallel, open-vocabulary detection and phrase grounding methods typically rely on large vision--language pretraining and retain specialized decoding for boxes/masks, \eg, OWLv2~\citep{minderer2024scalingopenvocabularyobjectdetection}, GroundingDINO~\citep{Liu2023GroundingDM}, GLIP~\citep{li2022grounded}, and MDETR~\citep{kamath2021mdetr}. Falcon Perception targets a simpler interface: early fusion of image and text tokens using a shared backbone over a unified sequence modeling objective.

\noindent \textbf{Autoregressive token interfaces for perception:} Pix2Seq~\citep{chen2022pix2seq} shows that detection can be cast as language modeling by discretizing structured outputs and training with maximum likelihood. This provides a clean ``token interface'' to perception, but also highlights a main limitation: autoregressive decoding becomes expensive as output sequences grow, and dense prediction must avoid duplication and drift. Falcon Perception inherits the token-interface idea but is designed around dense instance segmentation solving the efficiency issue by using specialized heads. UViM~\citep{kolesnikov2022uvim} targets unifying multiple dense vision tasks and explicitly addresses the cost of modeling high-dimensional outputs by using a learned guiding code. In parallel, recent VLMs can output boxes/masks or interleave language with masks (\eg, LISA~\citep{lai2024lisa}, GLaMM~\citep{rasheed2024glamm}) and generalist systems (\eg, Gemini~\citep{comanici2025gemini}, Molmo~\citep{molmo}, Qwen3-VL series~\citep{qwen3}) demonstrate broad capabilities. Compact VLM releases such as Moondream2~\citep{moondream2} and Moondream3~\citep{moondream3} focus on: small, deployment friendly vision-language models. Falcon Perception is positioned as a perception first early-fusion design: a single dense stack optimized for dense grounding, with competitive results on segmentation benchmarks and strong OCR performance at small scale.

\vspace{0.5em}
\noindent \textbf{Scope:} Our goal is not to build a general-purpose VLM for open-ended reasoning (\eg, VQA with multi-step logic or long-form captioning). We focus on dense grounding regimes where the primary difficulty is \emph{localization} under open vocabulary: selecting the correct instance given a phrase (attributes, OCR text, etc) and producing an accurate mask. In this setting, early fusion is particularly natural: the model benefits from allowing prompt to condition on the visual tokens and task context throughout the stack, instead of restricting language-vision interaction to a late cross-attention module.
